% D-glossary.tex: terms, defined the way the paper uses them.
\section{Glossary}
\label{app:glossary}

\begin{description}
  \item[Anchor set.] The published issuer keys a relying party has chosen to trust. Trust in Polaris
    is always membership in the verifier's own anchor set, never a decision the issuer makes for it.
  \item[Anonymity set.] The crowd a membership proof hides in: the members of the epoch it was made
    against. Nothing in the schema floors it above one, and no verdict reports it.
  \item[Append-only.] A table on which UPDATE and DELETE are refused by a database trigger, so history
    can only be added to. The audit of record is the named pattern of applying this per table.
  \item[Assurance level.] What an enrolment can honestly claim about how the person was proven, derived
    from the evidence recorded beside it and never typed by an operator.
  \item[Attestation.] A directional, per-context statement that one agency accepts another's
    credentials, recorded as a row. Never transitive.
  \item[Authenticity pack.] A credential as a file: the token value, its signature and the issuer's
    public key. What a holder holds and a stranger verifies.
  \item[Canonicalisation.] Turning a structure into one exact sequence of bytes (sorted keys, no
    whitespace) so that a signature made on one machine can be checked on another that rebuilds the
    same bytes.
  \item[Consistency proof.] A short proof that a smaller Merkle tree is a prefix of a larger one, so
    that a log only appended between two heads.
  \item[Detached verifier.] A program that checks Polaris artifacts with no Polaris code, no database
    and no network: the \code{polaris-verify} package, and the script that is now a shim onto it.
  \item[Disclosure level.] What a verification records: zero-knowledge (no token identifier),
    selective (named attributes), or full (the identity, logged).
  \item[Duress-aware.] The word the front door uses for the duress code, because the mechanism
    notices duress and does not resist compulsion; the stronger word is refused by a check.
  \item[Epoch.] A closed, immutable commitment to the active-token set at a moment, as a Merkle root,
    against which membership proofs are made.
  \item[Equivocation.] A log telling two observers two different histories. The proof of it is two
    signed heads of the same size with different roots.
  \item[External noun.] A named party outside the repository that did something to it, on a date,
    with a result that can be pointed at: the only kind of entry the scoreboard accepts.
  \item[Federation.] Independent authorities, each its own root, recognising one another through
    explicit attestations, with no central service.
  \item[Gossip.] Witnesses sharing the log heads they were shown so that a split view is detected by
    comparison.
  \item[High availability.] A database that survives the loss of its leader by promoting a replica
    under a lease, drilled and measured.
  \item[Inclusion proof.] A short proof that one entry is in a Merkle tree with a given root.
  \item[Long-term validation.] Deciding that a signature was valid at the instant it was made, from
    evidence fixed at that instant, so that the decision survives the key's later retirement.
  \item[Merkle tree.] A cryptographic fingerprint of a whole set, built by hashing pairs upward to one
    root, that lets a small proof show one item belongs to the set.
  \item[ML-DSA.] The lattice-based signature scheme standardised as FIPS 204; parameter sets 65 and
    87 are accepted. Post-quantum in its assumptions; the repository claims agility, not settled
    security.
  \item[Mutation drill.] A program that removes or inverts a guarantee, a constraint, a trigger, a
    refusal, a check's inputs, and requires the tests that supposedly prove it to go red, with a
    negative control so that zero survivors is distinguishable from a harness that ran nothing.
  \item[Nullifier.] A value derived from a secret and a scope that lets a verifier detect a second
    use without linking uses across scopes. In Polaris \code{Poseidon(secret\,||\,scope\,||\,epoch)},
    a public input of the membership proof, per relying party and per epoch, never global.
  \item[OpenID4VP.] OpenID for Verifiable Presentations, the protocol by which a wallet presents a
    credential to a verifier; the path on which the first outside wallet spoke to Polaris.
  \item[Operating contract.] The rule adopted on 13 September 2026 that the unit of progress is a
    surviving external dependency, that a product change needs a named outside reason, and that
    everything else goes to the lab.
  \item[Pairwise identifier.] A different handle for each relying party, so that two parties cannot
    join their logs. In Polaris the login token's subject and the presentation handle are derived
    under the relying party's scope; it bounds what a verifier stores, not what it is shown.
  \item[Partition.] A table split into monthly pieces by timestamp, transparent to queries, so that
    old months can be detached to the archive.
  \item[PKCE.] A proof-of-possession extension to the authorization-code login flow, so that an
    intercepted code cannot be exchanged by an attacker.
  \item[Positive control.] The half of a measurement that proves the instrument can see anything: an
    adversary that must succeed on a linkable population, a mutated tree the check must fail on, a
    patched verifier the suite must notice.
  \item[Post-quantum.] Cryptography whose security does not rest on problems a quantum computer
    solves efficiently. In this paper, a statement about algorithms, never about the system.
  \item[Release candidate.] The version \code{1.0.0-rc.7}: every condition the operating contract
    names for 1.0.0 holds, and no operator other than the author has run it.
  \item[Replica.] A copy of the database that follows the primary with some lag; safe for reading
    immutable material, unsafe for deciding current authorization.
  \item[Revocation feed.] A signed list of the hashes of revoked credentials an authority issued,
    consumed offline.
  \item[Scoreboard.] \code{lab/EXTERNAL-NOUNS.md}: the file that decides whether the project
    continues, holding only what named outside parties did, with zero and blank as valid values.
  \item[SD-JWT VC.] A credential format in which the issuer signs digests of the claims and the holder
    discloses the ones a verifier asked for; the format the first outside wallet presented.
  \item[Signed tree head.] A log's signed commitment to its whole history at a size.
  \item[Standby region.] A second cluster with its own lease store, streaming asynchronously from the
    first, electing no primary of its own until promoted; its recovery point is measured, not zero.
  \item[Status assertion.] A short-lived, issuer-signed statement of a credential's current status,
    stapled to a presentation so that authorization is decidable offline.
  \item[Sunset.] The day the old credential stops being accepted, treated as a decision with a floor
    under it because it is the moment the new one becomes compulsory.
  \item[Trust list.] A signed statement of every key an instance knows, with each key's status and
    the instants those statuses took effect.
  \item[Trusted referee.] A person whose vouching lets somebody with no documents be enrolled, bounded
    and co-signed past a threshold, and invisible on the credential.
  \item[Two-witness principle.] No cryptographic verdict is trusted from a single implementation; a
    second, independently written one must agree or abstain explicitly.
  \item[Witness.] An independent party that cosigns log heads it found consistent and refuses a fork.
  \item[Zero-knowledge proof.] A proof that a statement is true which reveals nothing beyond the
    statement's truth; in Polaris, membership of a token in an epoch.
\end{description}

\subsection*{Abbreviations}

Every abbreviation the text uses, spelled out, with what each standard or body is.

\begin{description}
  \item[API.] Application programming interface: the requests one program uses to call another.
  \item[CBOR.] Concise Binary Object Representation, the binary data format ISO 18013-5 mobile documents use.
  \item[CI.] Continuous integration: the build, tests and drills run automatically on every push.
  \item[CSRF.] Cross-site request forgery: an attack in which another site makes a browser send a request as a signed-in user.
  \item[ES256.] The classical ECDSA signature over curve P-256 that the HAIP profile pins.
  \item[FIPS.] US Federal Information Processing Standards; FIPS 202 defines SHA-3, FIPS 204 ML-DSA and FIPS 205 SLH-DSA.
  \item[HA.] High availability.
  \item[HAIP.] The High Assurance Interoperability Profile of OpenID4VP, which pins the credential format, the signature algorithm and the response mode a wallet and a verifier must use.
  \item[HTTP.] Hypertext Transfer Protocol, over which web services exchange requests.
  \item[ISO/IEC 18013-5.] The international standard for the mobile driving licence.
  \item[JSON.] JavaScript Object Notation, the text format every signed Polaris statement is written in.
  \item[KMS.] Key management service: a cloud service that holds keys.
  \item[NFC.] Near-field communication: the short-range radio a card answers a reader over.
  \item[NIST SP 800-63.] The US National Institute of Standards and Technology's digital identity guidelines, which define assurance levels.
  \item[npm.] The public package registry for JavaScript and TypeScript.
  \item[PIN, PUK.] The personal code that unlocks a card, and the unblocking key that clears a card locked by too many wrong PINs.
  \item[PKCS\#11.] The standard programming interface to hardware security modules and tokens that generate and hold keys inside themselves.
  \item[PyPI.] The Python Package Index, the public registry for Python packages.
  \item[QR.] A two-dimensional barcode read by a camera; one way a presentation is transferred.
  \item[RFC.] Request for Comments: the series in which internet standards are published; RFC 6962 describes Certificate Transparency.
  \item[SDK.] Software development kit: a library a relying party builds into its own program.
  \item[SHA3-256.] A cryptographic hash function of the SHA-3 family (FIPS 202) giving a 256-bit digest; Polaris signs the SHA3-256 digest of a token value.
  \item[SLH-DSA.] The hash-based signature scheme standardised as FIPS 205, registered as the fallback should the lattice assumptions fall.
  \item[SQL.] The query language of relational databases, in which the Polaris schema, triggers and procedures are written.
  \item[TLA+.] A formal specification language whose model checker explores every interleaving of events to show a property holds.
  \item[TLS.] Transport Layer Security, the protocol of the encrypted connection; the Polaris edge negotiates a hybrid post-quantum key exchange.
  \item[W3C.] The World Wide Web Consortium, which publishes the Verifiable Credentials data model.
  \item[WebAuthn.] The standard for proving possession of a hardware key or a device's built-in authenticator; operators use it as a second factor after the password.
  \item[X.509.] The standard format of a public-key certificate.
  \item[ZK.] Zero knowledge.
\end{description}
